%=====================================================================
\section{Every code artifact}
\label{sec:artifacts}

One card per file: its role, where it sits (Fig.~\ref{fig:map}), what it takes and gives, the logic it follows,
why it is needed, the alternatives, and what checks it. Line counts are those of the current version.

%---------------------------------------------------------------------
\subsection{Data plane (P4)}

\begin{card}{\code{p4/headers.p4} (323 lines)}
Role & every header and metadata type: Ethernet, up to two VLAN tags, IPv4 (+options), IPv6, TCP, TCP option bytes
(stack of 40), UDP; \code{packet\_in} (offset, fromClient, msgStart, msgType, connId) and \code{packet\_out} (kind, index); the metadata structs
of each walk (\code{tls\_meta\_t}, \code{http\_meta\_t}, \code{dhcp\_meta\_t}, \code{x509\_meta\_t}, TCP options);
\code{hs\_digest\_t}; the capacities; \code{CONN\_SLOTS} = 65,536, \code{COPY\_SESSION} = 100, \code{CPU\_PORT} = 255. \\
Logic & metadata fields that must survive a resubmit or a clone are annotated \code{@field\_list(0)} (arrival port,
pending handshake record) or \code{@field\_list(0,1)} (the copy's offset, direction, start, type, and the
connection's number). \\
Feeds & every other P4 file (included); P4Info names used by the controller. \\
Why & one place for every layout: a digest struct here and the controller's field list must agree in order. \\
\end{card}

\begin{card}{\code{p4/parser.p4} (237 lines)}
Role & \code{MyParser}: Ethernet, up to two VLAN tags (802.1Q; 802.1ad outside 802.1Q), IPv4/IPv6, TCP and its
options, UDP; includes the four payload walks. \\
Inputs & the packet; \code{instance\_type} and \code{ingress\_port} (resubmit $\to$ Ethernet only; port 255 $\to$
\code{packet\_out} first). \\
Outputs $\to$ & headers and metadata to ingress: addresses (as 128 bits for both versions), TTL, payload length,
option kinds (shifted into a 320-bit field), MSS, window scale. \\
Logic & option walk: \code{left} = data offset $\times$ 4 $-$ 20; each byte extracted into the stack (so the deparser
re-emits it), values read with \code{lookahead} before extraction; EOL and NOP one byte, the walk continues after
EOL; malformed length $\to$ the rest kept as padding (Fig.~\ref{fig:opt}). Only the first IPv4 fragment is read.
VLAN tags are extracted and re-emitted by the deparser; a third tag: forwarded, not fingerprinted. \\
Why & JA4T/JA4TS need every option kind in order; MSS and window scale values. \\
Checked by & JA4T 62/62, JA4TS 51/51 exact; \code{test\_vlan\_and\_ssh\_port} (one and two tags). \\
\end{card}

\begin{card}{\code{p4/parser\_tls.p4} (356 lines)}
Role & ClientHello / ServerHello walk for JA4 / JA4S; also decides ``incomplete'' and short segments. \\
Inputs & TCP payload; \code{payLen}; packet-out kind (3 $\to$ certificate walk). \\
Outputs $\to$ & \code{meta.tls}: hsType, versions, SNI flag, cipher(s), extension list, signature algorithms,
first ALPN name (length, first and last byte), overflow bits, \code{walked} / \code{parsed} / \code{incomplete} $\to$
Stream, Tls, main (resubmit). \\
Logic & 9-byte peek (record + handshake header); complete if hsLen + 9 $\le$ payLen and hsLen + 4 $\le$ recLen
(the record may continue: ServerHello and Certificate often share one record); each list read value by value with
\code{lookahead}, skipped with \code{advance}; GREASE skipped in the \code{select}; past the capacity: overflow bit.
Segments of 1--8 bytes starting with 0x16 or a method's first letter are treated as possible message starts.
Its first state also notes whether the payload begins with ``SSH-'' (\code{meta.sshBanner}, for Ssh). \\
Why & every JA4 / JA4S input comes from these lists. \\
Alternatives & see D2, D11, D12. \\
Checked by & JA4 152/152, JA4S 92 + 1 explained; split/reorder 202/202; capacity test. \\
\end{card}

\begin{card}{\code{p4/parser\_http.p4} (395 lines)}
Role & HTTP/1.x request walk for JA4H. \\
Inputs & a payload starting with one of nine methods (4-byte check). \\
Outputs $\to$ & \code{meta.http}: method bytes, version (last 8 bytes of the request line), header names chain
(512 B), header count, Cookie and Referer flags, cookie pairs chain (1,024 B), Accept-Language (4 chars), flags
$\to$ Http, Stream. \\
Logic & one byte per state (\code{HTTP\_READ\_BYTE}: lookahead, advance, left$-$1); names appended as read, with a
lower-cased copy of the last 16 bytes for matching (bit 5 set); Cookie / Referer names removed afterwards by
constant 6- / 7-byte slices; cookie value split at ``; '' with a zero byte between pairs; bare LF accepted; payload
ends before the empty line $\to$ \code{incomplete}. \\
Why & JA4H needs names in order, the cookie names and pairs, and the language. \\
Checked by & JA4H 72/72; split/reorder (HTTP cases); capacity test (names, cookies). \\
\end{card}

\begin{card}{\code{p4/parser\_dhcp.p4} (354 lines)}
Role & DHCPv4 (67/68) and DHCPv6 (546/547) option walks for JA4D / JA4D6. \\
Logic & v4: skip the 240-byte BOOTP header and magic cookie; per option: code joins the list except 0, 50, 53, 81,
255; 53 type, 57 max size, 50 requested IP, 81 domain if the name is non-empty, 55 request list. v6: type, 3-byte
transaction id, then 2+2-byte options; containers (IA\_NA 3, IA\_PD 25: 12 bytes; IAADDR 5: 24; IAPREFIX 26: 25;
IA\_TA 4: 4 and the \code{i} flag) skip only their fixed fields so nested options come next in packet order; 1 gives
the DUID length; 39 sets \code{d} only for a non-empty name (fixed after the live run); 6 the requested codes. \\
Outputs $\to$ & \code{meta.dhcp} $\to$ Dhcp; \code{walked} $\to$ resubmit. \\
Checked by & JA4D/D6 10/10 on FoxIO's samples; live run; \code{test\_dhcp6\_fqdn}, capacity test. \\
\end{card}

\begin{card}{\code{p4/parser\_x509.p4} (371 lines)}
Role & DER walk of one certificate (packet-out kind 3) for JA4X. \\
Logic & one shared \code{tlv} reader (tag, short or long-form length) entered with a return code \code{ret}
(\code{R\_CERT} \ldots\ \code{R\_EXT}); countdowns per level (tbs, name, set, item, extensions) mark where each
SEQUENCE / SET ends; each OID's content copied as $\langle$len$\rangle\langle$bytes$\rangle$ into the issuer,
subject or extension list (256 B each). \\
Outputs $\to$ & \code{meta.x509} $\to$ X509. \\
Checked by & JA4X 26/26; live run (chain of 2). \\
\end{card}

\begin{card}{\code{p4/main.p4} (213 lines)}
Role & ingress, egress, deparser; ConnKey; ConnId (the connection's number, D21); the \code{forward} table. \\
Logic & Fig.~\ref{fig:life} and the ingress flowchart: arrival port kept across resubmit; packets from the controller
dropped after reading; controls in a fixed order; resubmit after a walk; the handshake digest from one call,
postponed to pass 2 when this pass sent another digest (a message's, or Ssh's); egress E2E clone when
\code{copy}, \code{packet\_in} filled on the clone (with the connection's number); deparser emits \code{packet\_in},
Ethernet, VLAN tags, IP, TCP (+options), UDP. ConnId: a new number from a counter when the slot's check differs or
a SYN brings a new initial sequence number; packets from the controller keep the slot's number. \\
Why & the glue; every ordering rule here answers a target constraint (D5, D7, D9). \\
Checked by & forwarding 9,013/9,013 frames unchanged; every fingerprint test; \code{test\_port\_reuse}. \\
\end{card}

\begin{card}{\code{p4/handshake.p4} (127 lines)}
Role & JA4T, JA4TS, JA4L: state machine per slot (EMPTY, SYN, SYNACK, DONE, SYNACK\_ONLY). \\
Logic & first SYN: store time A, client TTL, report JA4T; first SYN-ACK after it: time B, server TTL, JA4TS; the
first ACK without SYN after that: latencyC = (C$-$B)/2, latencyS = (B$-$A)/2, JA4L; a SYN-ACK without its SYN: JA4TS
only. The record goes to \code{meta.hsd} (sent by main.p4). \\
Checked by & JA4T/TS exact; TTL 102/102; \code{test\_ja4l\_gaps}. \\
\end{card}

\begin{card}{\code{p4/stream.p4} (209 lines)}
Role & which segments the controller must see (copy rules (a)--(e), Fig.~\ref{fig:copyrules}). \\
State & per slot: connection number, client side, ISNs, started, pending, pendingSeq, edge, gap start, gap open,
HTTP seen (per direction). \\
Outputs $\to$ & \code{meta.copy}, offset, direction, msgStart, msgType $\to$ egress clone; \code{fromClient}. \\
Logic & offset = seq $-$ ISN $-$ 1; window 16 KB; a returned joined message (or the certificates) clears pending; a
retransmission bringing the whole message also clears it. Rule (e): the edge follows the furthest byte; a segment
beyond it opens a gap (in a direction that carried an HTTP request, copied within 16 KB); a segment covering the
gap's first byte closes it. Sequence numbers are compared modulo $2^{32}$ (signed difference). \\
Why & the switch must copy every piece a split message may need, and nothing else. \\
Checked by & 202/202 split/reorder cases; 8 limit cases never wrong; \code{test\_keepalive\_reorder} 18/18;
\code{test\_server\_flight\_reorder} 19/19. \\
\end{card}

\begin{card}{\code{p4/tls.p4}, \code{http.p4}, \code{ssh.p4}, \code{dhcp.p4}, \code{x509.p4}}
Role & one control per fingerprint family: decide when to report, keep the per-connection state for that, send one
digest (Section~\ref{sec:channels}). \\
Logic & Tls: once per direction (\code{tlsDone} bits, retransmissions not twice). Http: one per request, not twice
for one sequence number; the request's sequence number in the digest. Ssh: a connection on port 22 or with an
``SSH-'' banner (bare ACKs counted from the start); counts, batches of 8, window of 200, flush; on a reused port's
SYN the previous connection's last window. Dhcp: every message. X509: every certificate with its index. DHCP, HTTP
and SSH digests carry a serial number (D10); every TCP digest the connection's number (D21). \\
\end{card}

%---------------------------------------------------------------------
\subsection{Controller (Python)}

\begin{card}{\code{controller/p4runtime\_client.py} (201 lines)}
Role & a minimal P4Runtime client: arbitration (become primary), \code{SetForwardingPipelineConfig}, table writes,
clone session, digest configuration and acknowledgement, packet-in callback, packet-out with metadata. \\
Logic & a background thread reads the stream; each digest list goes to its handler, then is acknowledged; P4Info is
used to translate names to ids. \\
Why & the only path between the controller and the switch. \\
Alternatives & \code{p4runtime-shell}: interactive, heavier; this client does only what is needed and is easy to
read. \\
\end{card}

\begin{card}{\code{controller/main.py} (240 lines)}
Role & the Controller: start-up (program, forwarding, clone session, six digests), one handler per digest, packet-in
to the reassembler and packet-out back, \code{finish()} (SSH flush). \\
Outputs $\to$ & records (\code{records.py}), console lines, JSON file. \\
Logic & each handler decodes its struct (field order = the P4 struct), finds the connection's record by its
endpoints and number, sets the client / server roles (the SYN sender, the ClientHello sender, the SSH digest's
direction), writes the text via the formatting module. JA4H are kept in stream order by each request's sequence
number. An SSH connection is remembered for the final flush; an earlier one on the same endpoints is not flushed
(its window came on the new SYN). \\
\end{card}

\begin{card}{\code{controller/handshake.py}, \code{tls.py}, \code{http.py}, \code{ssh.py}, \code{dhcp.py}, \code{x509.py}}
Role & decode one digest struct and write FoxIO's text; no packet bytes. \\
Logic & handshake: \code{window\_kinds\_mss\_scale}, JA4L halves. tls: version code (supported\_versions if present),
SNI flag, two-digit counts (capped at 99 as FoxIO), ALPN first / last characters ('9' for a non-ASCII byte, '0' for a missing one, \code{00} without ALPN), sorted
ciphers; sorted extensions without SNI and ALPN; signature algorithms in order; SHA-256 first 12 hex digits. http:
method code, version, flags, count, language; three hashes. ssh: histogram and mode per direction (smallest length on
a tie, as FoxIO). dhcp: type names, size, flags, lists. x509: OIDs as hex joined by ','. \\
Checked by & every EXACT row of the comparison. \\
\end{card}

\begin{card}{\code{controller/reassembly.py} (199 lines)}
Role & join split messages by framing (Section~\ref{sec:seq} and the algorithm in Section~\ref{sec:algos}). \\
Inputs & copied segments with (endpoints, connection number, direction), seq, offset, start flag and type. \\
Outputs $\to$ & complete messages with the sequence number of their first byte $\to$ packet-out kind 1, or a list of
certificates $\to$ packet-out kind 3 each. \\
Logic & bounded (D24): 32 KB per direction (lowest offsets dropped first), idle directions dropped after 60 s, at
most 100,000 directions; a running byte count, so the bound costs nothing per piece. \\
Checked by & split/reorder 202/202; \code{test\_controller\_state} (bounds, a false start inside a request). \\
\end{card}

\begin{card}{\code{controller/records.py} (57 lines)}
Role & one record per connection (key: endpoints in a fixed order and the switch's connection number), first value of
each name kept; lists for JA4H, JA4SSH, JA4D; per-message list for DHCP in arrival order. Records are kept until
the run ends and written to JSON (a long-running deployment would export and drop them; see the assumptions). \\
\end{card}

%---------------------------------------------------------------------
\subsection{Tools, tests, network}

\begin{card}{\code{tools/replay.py} (242 lines) and \code{tools/nanoport.py} (57 lines)}
Role & run the real switch and controller without root and send a capture through them. \\
Logic & free gRPC / Thrift ports; ipc path in the runtime folder; switch dies with its parent; one nanomsg socket,
opened after the switch answers; warm-up frame; client / server side by the SYN sender; gaps kept (max 2 s);
\code{sent.pcap} with send times; frames out collected; forwarding count. \\
Outputs $\to$ & records, messages, \code{sent.pcap}, forwarding numbers $\to$ \code{compare.py} and the tests. \\
\end{card}

\begin{card}{\code{tools/reference.py} (172 lines)}
Role & the references: FoxIO's Rust tool (JSON per connection), tshark fields for JA4TS, FoxIO's stored Wireshark
plugin output for JA4D, and the plugin's rule on tshark's PDML for captures without stored output (skipping DHCP
quoted inside ICMP errors). A later TCP connection on the same endpoints is keyed with its place among them, from
tshark's stream numbers (which FoxIO's tool reports). \\
\end{card}

\begin{card}{\code{tools/compare.py} (219 lines)}
Role & per capture: replay, then a verdict per fingerprint per connection; totals; forwarding; JA4L distribution;
not-compared counts (HTTP/2, TLS 1.3 certificates). \code{compare\_records} is shared with the live run. \\
Outputs $\to$ & console report, \code{results/compare\_all.json} $\to$ the results tables of both PDFs. \\
Logic & EXPLAINED only for differences checked by hand and written with their reason; EXCLUDED only tunnel captures;
exit code non-zero on any WRONG or any altered frame. The switch's records are keyed like the references: endpoints,
and the place among connections on them in the order of the connection numbers. \\
\end{card}

\begin{card}{\code{tools/measure\_field\_counts.py} (326 lines)}
Role & the measurement study of list lengths (capacities). \\
Logic & local listener records the first message of each installed client; corpus of 42 captures; tshark PDML
(reassembled) for TLS / HTTP, the JA4D rule for DHCP, \code{cryptography} for certificates; per list: p50, p95, p99,
max, distinct, share within the capacity. \\
Outputs $\to$ & \code{results/measure/field\_counts.json} $\to$ \code{results\_field\_counts.tex}. \\
\end{card}

\begin{card}{\code{network/run\_live.py} (249 lines)}
Role & the live run: Mininet h1 -- s1 -- h2 with BMv2 on real interfaces, real servers and clients, captures on the
switch ports, comparison and forwarding check. \\
Logic & offloads off; certificates made for a chain of 2; servers' process ids kept; clients start once h2 listens
on every server port (\code{ss}); every client under \code{timeout -k 2 15}, its exit status checked; a Ctrl-C interrupts the running host command before cleanup; results chowned to the user. \\
Outputs $\to$ & \code{results/live/} (report, rows, records, captures) $\to$ \code{tools/live\_tex.py}. \\
\end{card}

\begin{card}{\code{tools/live\_tex.py} (45 lines)}
Role & writes the live run's results table for both PDFs, so the numbers are never typed by hand. \\
Logic & counts verdicts per fingerprint in \code{rows.json}; forwarding lines from \code{report.txt}; JA4L
latency differences recomputed from the rows (median as the live run prints it, the upper middle value). \\
Outputs $\to$ & \code{docs/results\_live.tex} $\to$ \code{design.pdf} and this document (\code{make live} runs it). \\
\end{card}

\begin{card}{\code{tests/} (11 files, 27 tests) and \code{Makefile}}
Role & tests in Section~\ref{sec:verif}; \code{make build / compare / test / live / docs / clean}. \\
\end{card}
